\chapter{Background \& Literature Review}
\label{ch:background}

This chapter reviews the ideas this dissertation is built on. It begins
with what reinforcement learning is and how it differs from the more
familiar kinds of machine learning; formalises the Markov Decision
Process, the mathematical object every later chapter manipulates; then
traces the two algorithm families compared in this project --- value-based
methods (Q-learning to DQN) and policy-gradient methods (REINFORCE to
PPO) --- from their origins to their modern forms. The second half turns
to the application domain: what the literature reports about
reinforcement learning for trading, the recurring evaluation mistakes
that make much of that literature unreliable, and the partial-observability
question that markets pose. The chapter closes by stating the gap this
project occupies.

\section{From trial and error to a formal framework}
\label{sec:bg_intro}

Most machine learning taught at masters level is \emph{supervised}: a
model is shown inputs together with correct answers (this image is a cat;
this email is spam) and learns to reproduce the labelling. Reinforcement
learning \cite{sutton2018} removes the answer sheet. The learner --- the
\emph{agent} --- interacts with a world --- the \emph{environment} --- in
a loop: it observes the situation, chooses an action, and receives back
two things only: a new situation and a numerical \emph{reward}. Nobody
ever tells it which action was correct; it must discover, by trial and
error, which behaviour accumulates the most reward over time.

Two consequences of removing the answer sheet drive everything in this
dissertation. First, \emph{credit assignment}: when reward arrives late
(the month's profit is only known at the month's end), the agent must
work out which of its many earlier decisions deserve the credit or the
blame. Second, \emph{verification}: with no labels, there is no
training-set accuracy to inspect, so a defective implementation produces
no error --- only quietly wrong behaviour. The first consequence
motivates the mathematics of Chapter~\ref{ch:methodology}; the second
motivates this project's game-first validation discipline
(Chapter~\ref{ch:implementation}), because games are the rare setting
where the correct outcome of learning \emph{is} known.

The field's landmark systems illustrate the paradigm's reach:
TD-Gammon learned expert backgammon purely from self-play
\cite{tesauro1995}; DQN reached human-level play across dozens of Atari
games from raw pixels \cite{mnih2015}; AlphaGo defeated a world champion
at Go \cite{silver2016}. All three learned from a scalar score alone ---
precisely the position a trading agent is in, where the score is money.

\section{The Markov Decision Process}
\label{sec:bg_mdp}

The standard formalism for sequential decision problems is the Markov
Decision Process (MDP), whose foundations trace to Bellman's dynamic
programming \cite{bellman1957}; the modern treatment followed throughout
this dissertation is Sutton and Barto \cite{sutton2018}. An MDP is a
tuple
\[
(\mathcal{S}, \mathcal{A}, P, R, \gamma),
\]
where $\mathcal{S}$ is the set of \emph{states} the world can be in,
$\mathcal{A}$ the set of \emph{actions} available, $P(s' \mid s, a)$ the
\emph{transition dynamics} (the probability of landing in state $s'$
after taking action $a$ in state $s$), $R$ the \emph{reward function},
and $\gamma \in (0,1]$ a \emph{discount factor} that down-weights distant
rewards. The agent's behaviour is a \emph{policy} $\pi(a \mid s)$: a rule
(possibly probabilistic) for choosing actions in each state. The learning
goal is to find a policy maximising the expected \emph{return} --- the
cumulative discounted reward
$G_t = \sum_{k \ge 0} \gamma^k r_{t+k}$.

\paragraph{The Markov property.} The definition requires that the state
summarise \emph{all decision-relevant history}: given the present state,
the past adds nothing. In plain terms, a chess board satisfies this ---
one photograph of the board is enough to play on; how the pieces got
there is irrelevant. A poorly chosen trading state violates it: if the
agent is shown only today's price change but its profit depends on the
price it originally paid, then something decision-relevant is hidden, and
the theory's guarantees quietly lapse. This single observation drives the
central design decision of this project --- the \emph{complete} state of
Section~\ref{sec:state} --- and its most important limitation analysis
(Section~\ref{sec:bg_pomdp}).

\paragraph{Value functions and the Bellman equation.} Two functions
organise the solution space. The state-value function
$V^\pi(s) = \mathbb{E}_\pi[\, G_t \mid s_t = s\,]$ scores how good a
situation is under policy $\pi$; the action-value function
$Q^\pi(s,a) = \mathbb{E}_\pi[\, G_t \mid s_t = s, a_t = a\,]$ scores how
good an action is in that situation. The optimal action-value function
obeys Bellman's recursion,
\begin{equation}
Q^\ast(s,a) \;=\; \mathbb{E}\!\left[\, r + \gamma \max_{a'} Q^\ast(s', a')
\;\middle|\; s, a \right],
\label{eq:bg_bellman}
\end{equation}
which reads: the value of acting now is the immediate reward plus the
discounted value of acting optimally afterwards. Equation
(\ref{eq:bg_bellman}) is not merely notation --- it is the operating
principle of the entire value-based family reviewed next, and its
sampled, approximate form is the loss function our Deep Q implementation
minimises (Section~\ref{sec:dqn}).

\paragraph{Exploration versus exploitation.} A learner that always takes
the action it currently believes best may never discover a better one; a
learner that experiments forever never banks its knowledge. Every
algorithm in this dissertation embodies a stance on this dilemma:
value-based methods explore explicitly (taking a random action with
probability $\varepsilon$), while policy-gradient methods explore
implicitly, by sampling actions from a probabilistic policy that only
gradually sharpens.

\section{Value-based methods: from Q-learning to DQN}
\label{sec:bg_value}

The value-based family answers ``what should I do?'' indirectly: learn
how much each action is \emph{worth}, then pick the most valuable. Its
lineage begins with temporal-difference (TD) learning \cite{sutton1988},
which updates value estimates from the difference between successive
predictions rather than waiting for final outcomes, and reaches its
canonical form in Q-learning \cite{watkins1992}: after each transition
$(s, a, r, s')$, update
\begin{equation}
Q(s,a) \;\leftarrow\; Q(s,a) + \alpha \Bigl[\,
r + \gamma \max_{a'} Q(s', a') - Q(s,a) \Bigr],
\label{eq:bg_qlearning}
\end{equation}
nudging the stored value toward the Bellman target with learning rate
$\alpha$. Remarkably, this converges to $Q^\ast$ for finite MDPs
regardless of how actions were chosen during learning (it is
\emph{off-policy}), provided every state--action pair is visited
sufficiently often \cite{watkins1992}.

\paragraph{The table problem.} Tabular Q-learning stores one number per
state--action pair, which is fine for a small grid-world and hopeless for
anything continuous. An intuition used repeatedly in this project's
supervision: imagine a haulage robot whose state is its position in a
mine. Discretise the mine coarsely and the table is small but the robot
is blind to distinctions that matter; discretise finely and the table
outgrows memory --- and either way, experience at one location teaches
nothing about the location one metre away, because table cells share
nothing. Our own Phase-0 experiments reproduce this on CartPole: tabular
Q-learning succeeds \emph{only after} the 4-dimensional continuous state
is hand-binned, and the approach is a non-starter for the 9-dimensional
continuous trading state (Section~\ref{sec:res_phase0}).

\paragraph{Function approximation and its dangers.} The remedy is to
replace the table with a function approximator --- a neural network
$Q_\phi(s,a)$ --- that \emph{generalises} across nearby states. The cost
is that convergence guarantees weaken substantially: the combination of
function approximation, bootstrapping and off-policy data (the ``deadly
triad'') can diverge \cite{tsitsiklis1997, sutton2018}. Deep Q-learning
(DQN) \cite{mnih2015} demonstrated that two mechanisms restore practical
stability: \emph{experience replay} \cite{lin1992} --- storing
transitions in a buffer and training on random minibatches, which breaks
the strong temporal correlation of consecutive experiences --- and a
periodically frozen \emph{target network}, which stops the regression
target from chasing its own updates. Both mechanisms are implemented from
scratch in this project (Section~\ref{sec:impl_dqn}), and both matter
doubly for market data, whose consecutive samples are even more
correlated than game frames. Refinements such as double Q-learning
\cite{vanhasselt2016} address residual overestimation bias; they are
noted as extensions rather than implemented, to keep the from-scratch
implementation minimal and fully understood.

\section{Policy-gradient methods: REINFORCE and its descendants}
\label{sec:bg_pg}

The policy-gradient family answers ``what should I do?'' directly: write
the policy as a differentiable function $\pi_\theta(a \mid s)$ --- in
practice a neural network whose weights are $\theta$ --- and adjust
$\theta$ by gradient ascent on the expected return $J(\theta)$. The
foundational algorithm is REINFORCE \cite{williams1992}. Its enabling
mathematical step is the likelihood-ratio (score-function) identity
$\nabla_\theta p_\theta = p_\theta \nabla_\theta \log p_\theta$, which
converts the gradient of an expectation \emph{over} a distribution into
an expectation \emph{of} a gradient --- something that can be estimated
from sampled episodes. This identity is the sole reason a
$\log \pi_\theta$ term appears in every policy-gradient loss; the full
six-step derivation, each step labelled exact or approximate, is given in
Section~\ref{sec:pg_derivation}.

REINFORCE's virtue is transparency: it is unbiased and involves no
bootstrapping, no replay, and no auxiliary networks. Its vice is
\emph{variance}: each gradient estimate is scaled by a full Monte-Carlo
return, a noisy quantity, so learning is slow and unstable. The
subsequent thirty years of policy-gradient research can be read as a
sequence of patches on this one weakness:

\begin{itemize}
  \item \textbf{Baselines and advantage.} Subtracting a state-dependent
        baseline $b(s)$ from the return leaves the gradient unbiased but
        can reduce variance dramatically; learning that baseline with a
        second network (a \emph{critic}) gives actor-critic methods
        \cite{konda2000}, whose modern parallel form is A2C/A3C
        \cite{mnih2016}. Generalised advantage estimation
        \cite{schulman2016gae} tunes the bias--variance trade-off of the
        advantage signal itself.
  \item \textbf{Constrained updates.} Even with a good advantage signal,
        a single large policy update can be destructive, because the data
        was collected under the \emph{old} policy. Trust-region methods
        \cite{schulman2015trpo} bound the update size formally; Proximal
        Policy Optimization (PPO) \cite{schulman2017ppo} achieves a
        similar effect by simply \emph{clipping} the update ratio, and
        has become the default policy-gradient method in applied work,
        including most trading-RL frameworks \cite{liu2020finrl}.
\end{itemize}

This project retains REINFORCE as the primary policy-gradient method
precisely \emph{because} it is the un-obscured foundation: every design
element of the successors is a response to a weakness REINFORCE exhibits
transparently, and our Phase-0 results reproduce those weaknesses on cue
(CartPole: REINFORCE 292 vs.\ PPO 500,
Section~\ref{sec:res_phase0}). PPO is retained as the ablation's third
row so the patched descendant can be compared with its ancestor on the
same problem.

\section{Choosing between the families}
\label{sec:bg_families}

The two families make opposite engineering trade-offs, summarised in
Table~\ref{tab:bg_families}. Neither dominates: the value-based route
updates from every single step and reuses old experience, but bootstraps
(estimates built on estimates) and only \emph{derives} a policy
indirectly; the policy-gradient route optimises the actual objective with
unbiased gradients, but needs whole episodes per update and discards
experience once used. Comparing structurally different learners on the
\emph{same} problem is therefore scientifically informative --- if both
converge to the same behaviour, that behaviour likely reflects the
problem's true optimum rather than an algorithmic artefact. This logic is
the core experimental design of Chapter~\ref{ch:results}.

\begin{table}[h]
\centering
\caption{The two algorithm families compared in this project.}
\label{tab:bg_families}
\small
\begin{tabular}{lll}
\toprule
 & \textbf{Value-based (Deep Q)} & \textbf{Policy-gradient (REINFORCE)} \\
\midrule
Learns            & action values $Q_\phi(s,a)$ & the policy $\pi_\theta(a|s)$ itself \\
Update timing     & every environment step      & once per completed episode \\
Data reuse        & replay buffer (off-policy)  & fresh rollouts only (on-policy) \\
Estimator bias    & biased while learning (bootstraps) & unbiased \\
Estimator variance& low (one-step targets)      & high (full Monte-Carlo returns) \\
Exploration       & explicit ($\varepsilon$-greedy) & implicit (stochastic policy) \\
Failure style     & divergence/overestimation   & noisy, slow convergence \\
\bottomrule
\end{tabular}
\end{table}

\section{Reinforcement learning for trading}
\label{sec:bg_trading}

Applying RL to trading is an old idea with a well-documented set of
pitfalls. The foundational work is Moody and Saffell's \emph{direct
reinforcement} \cite{moody2001}, which trains a trading policy directly
by gradient methods and --- presciently for this project --- argues that
the \emph{objective} must be risk-adjusted (they propose a differential
Sharpe ratio), because maximising raw profit incentivises maximum
exposure. That is an early recognition of the theme our results confirm
empirically: \emph{the reward function defines the incentive}, and an
agent can only be blamed for optimising what it was actually asked to
optimise. Deng et al.\ \cite{deng2017} extended direct RL with deep
recurrent representations; Jiang et al.\ \cite{jiang2017} treat
multi-asset portfolio management with deep RL on cryptocurrency data.
Modern open frameworks, notably FinRL \cite{liu2020finrl}, standardise
Gym-style market environments and default to PPO-family agents ---
convenient, but, as this dissertation argues, convenience without
derivation invites exactly the silent failures that motivated our
from-scratch requirement.

\paragraph{A close reading of the nearest prior work.} The study closest
to ours is Th\'eate and Ernst \cite{theate2021}, who benchmark a trading
DQN (TDQN) against buy-and-hold and classic active strategies across 30
instruments. Read critically, their own results anticipate this
dissertation's central finding. Their headline is that TDQN
``outperforms benchmark strategies on average'' --- but the strategies it
beats on average are the \emph{active} ones (trend-following posts a
\emph{negative} average Sharpe on their testbench), while, in their own
words, TDQN ``only barely surpasses the buy and hold strategy on these
particular bullish markets.'' Most tellingly, on the S\&P~500 ETF ---
the very asset studied here --- their reported test Sharpe for TDQN is
0.834, \emph{identical} to buy-and-hold's 0.834; they observe that TDQN's
performance ``is identical or very close to the passive trading
strategies for multiple stocks'' because the agent ``efficiently learns
to tend toward a passive trading strategy.'' They further report
substantial run-to-run variance and overfitting on volatile single names
(Tesla). In short: a careful prior study, honestly read, already shows
value-based RL converging toward passivity on index products under
profit-seeking objectives --- a result their framing treats as a
footnote, and this dissertation places at the centre, explains
mechanistically, and reproduces with fully derived implementations.

\paragraph{Evaluation pitfalls.} Surveys of the area \cite{fischer2018}
and the financial machine-learning methodology literature
\cite{bailey2014, lopezdeprado2018} repeatedly flag failure modes that
this project's design guards against explicitly:

\begin{enumerate}
  \item \textbf{Weak baselines.} Much published work beats only strawman
        alternatives. The baseline that matters is buy-and-hold,
        \emph{properly capitalised} --- deploying the same capital the
        agent may deploy. Section~\ref{sec:defects} reports a live example
        from this project's own earlier code, where a baseline named
        buy-and-hold in fact deployed capital a tenth at a time and so
        understated the benchmark it stood for.
  \item \textbf{Look-ahead leakage.} Shuffling time-series data before
        splitting lets fragments of the future into training. All splits
        in this project are strictly chronological
        (Figure~\ref{fig:split}).
  \item \textbf{Ignored transaction costs.} Costless backtests reward
        hyperactive trading that fees would destroy. Here, fees live
        \emph{inside} the reward, so the agent pays them during learning,
        and Section~\ref{sec:res_fees} sweeps the assumed cost to test
        whether the agent's behaviour, not merely its profit, responds to
        it.
  \item \textbf{Backtest overfitting.} Repeatedly tuning on the same
        held-out period converts test data into training data
        \cite{bailey2014}. Our defence is procedural: fixed budgets and
        hyperparameters chosen on games, one evaluation protocol declared
        in advance, multiple seeds reported.
\end{enumerate}

\paragraph{The efficient-market null hypothesis.} Financial economics
supplies the prior against which any trading result must be read: in a
broadly efficient market \cite{fama1970}, persistent excess return over
buy-and-hold from public price history alone should not be expected, and
the risk--return trade-off \cite{markowitz1952} implies that reducing
risk, not only raising return, is a legitimate form of success (measured
here by the Sharpe ratio \cite{sharpe1994} and maximum drawdown). This
dissertation therefore treats ``the agent rediscovers buy-and-hold'' as a
scientifically meaningful finding rather than a failure --- and treats
any easy-looking victory as a signal to audit the experiment, not to
celebrate.

\section{Partial observability and hidden state}
\label{sec:bg_pomdp}

The MDP formalism assumes the agent sees a state satisfying the Markov
property. When decision-relevant structure is hidden --- as it plainly is
in markets, where order flow, news, and the strategies of other
participants never appear in a price series --- the problem is formally a
\emph{Partially Observable} MDP (POMDP) \cite{kaelbling1998}. The
principled treatment maintains a \emph{belief state}: a probability
distribution over what the hidden variables might be, updated as
observations arrive. The closely related hidden Markov model
\cite{rabiner1989} formalises the special case where an unobserved
discrete regime (e.g.\ ``calm market'' vs.\ ``turbulent market'')
generates the observations --- a natural reading of financial data, and
one raised explicitly in this project's supervision as the honest way to
think about what a price-only state misses.

Exact POMDP solutions are intractable for all but tiny problems
\cite{kaelbling1998}; practical systems either learn a memory (recurrent
policies) or --- the route taken here --- \emph{engineer the observation
to be as complete as practically possible} and state the residual gap
honestly. Chapter~\ref{ch:methodology} designs a state that includes
everything the agent can legitimately know (market summary statistics,
the episode clock, and its full balance sheet); Chapter~\ref{ch:conclusions}
returns to the POMDP view as the principal direction for future work.

\section{Summary and gap}
\label{sec:bg_gap}

The literature offers, on one side, two mature algorithm families whose
mechanics and trade-offs are well understood on games
(Sections~\ref{sec:bg_value}--\ref{sec:bg_families}); on the other, a
trading-RL corpus of genuine ideas \cite{moody2001, deng2017, theate2021}
interleaved with results weakened by the evaluation pitfalls of
Section~\ref{sec:bg_trading}. Between them sits a gap that is
methodological rather than algorithmic: \emph{few works formulate one
careful trading MDP and subject structurally different, fully understood
learners to an evaluation designed to make honest conclusions
detectable}. That is the gap this dissertation occupies: a single
formulation (Chapter~\ref{ch:methodology}); from-scratch implementations
whose every line is traceable to a derived equation
(Chapter~\ref{ch:implementation}, Appendix~\ref{app:trace}); validation
on known-outcome games first; and an evaluation against properly
capitalised baselines under chronological splits, fee-inclusive rewards
and multiple seeds (Chapter~\ref{ch:results}).

\paragraph{Chapter summary (plain language).} Reinforcement learning
means learning by trial and error from a numerical score. Two big
families exist: one learns \emph{how much each choice is worth} and picks
the most valuable; the other learns \emph{the choosing rule itself} and
nudges it toward whatever scored well. People have applied both to
trading for decades, but published results are often flattered by unfair
comparisons --- weak opponents, peeking at the future, ignoring trading
costs. Markets also hide information no price chart contains, which the
theory calls partial observability. This project's plan follows directly:
build one honest trading problem, implement both families from the ground
up, prove them on games first, and compare them fairly. The next chapter
builds that problem precisely.
